\begin{ouverture}[espace]{17}{Espace, pyramides et cônes}{Partie 4 · Géométrie et espace}
{Perspective · Patron · Repérage · Volume}
\objectif{Lire et produire des représentations de solides.}
\objectif{Se repérer par abscisse, ordonnée et altitude.}
\objectif{Calculer les volumes d’une pyramide et d’un cône.}
\prerequis{Aires, volumes du pavé et du cylindre, Pythagore.}
\end{ouverture}
\pagestyle{otMain}\essentiel\fontsize{11}{15}\selectfont
\phantomsection\label{cours-17-1}
\essSection{1}{Voir et représenter}
\begin{definition}Une pyramide a une base polygonale et des faces latérales triangulaires réunies en un sommet. Un cône de révolution est obtenu par la rotation d’un triangle rectangle autour d’un côté de l’angle droit. Sa base est un disque ; une génératrice relie le sommet au cercle de base.\end{definition}
\begin{savoirfaire}[Perspective et patron]La perspective cavalière indique les arêtes cachées en pointillés ; elle ne conserve pas toutes les longueurs ni tous les angles. Un patron est une figure plane qui permet de reconstituer le solide par pliage. Les longueurs qui se rejoignent doivent coïncider. Pour un cône, le bord courbe du secteur latéral a la longueur du cercle de base.\end{savoirfaire}
\label{fig-cours-17-1-cours-solides}\figureGeo{cours-solides}

\begin{definition}Dans un pavé muni d’un repère, un point a trois coordonnées $(x;y;z)$ : abscisse, ordonnée, altitude. Lire dans l’ordre des trois axes, pas dans l’ordre apparent des arêtes sur la feuille.\end{definition}
\clearpage
\phantomsection\label{cours-17-2}
\essSection{2}{Des volumes liés à une base}
\begin{propriete}Pour une pyramide, $V=\frac13\mathcal B h$, où $\mathcal B$ est l’aire de la base et $h$ la hauteur perpendiculaire à son plan. Pour un cône de rayon $r$, $V=\frac13\pi r^2h$. Un prisme ou cylindre de même base et de même hauteur a un volume trois fois plus grand.\end{propriete}
\begin{exemple}Une pyramide de base carrée de côté 6 cm et de hauteur 8 cm a pour volume $(1/3)\times36\times8=96$ cm$^3$. Un cône de rayon 3 cm et de hauteur 5 cm a pour volume $(1/3)\pi\times9\times5=15\pi$ cm$^3$.\end{exemple}
\begin{savoirfaire}[Choisir les bonnes données]Calculer l’aire de base avant le volume. La hauteur est perpendiculaire à la base : ce n’est ni une arête inclinée ni une génératrice. Si les unités diffèrent, convertir avant de les multiplier. Une aire en cm$^2$ multipliée par une hauteur en cm donne un volume en cm$^3$.\end{savoirfaire}
\begin{attention}La perspective peut écraser une hauteur ou un angle droit. Se fier aux données, aux codages et aux noms, sans mesurer le dessin. Pour un cône, le diamètre doit être divisé par 2 avant d’utiliser le rayon dans $\pi r^2$.\end{attention}

\clearpage\rubrique{Découverte}
\begin{blocLivre}{ex-17-1}\begin{exercice}[1]{Nommer les éléments}
Une pyramide SABCD a une base carrée ABCD. Nommer son sommet, sa base et le nombre de ses faces latérales.
\label{fig-ex-17-1}\figureGeo{sol-01}
\end{exercice}\end{blocLivre}
\begin{blocLivre}{ex-17-2}\begin{exercice}[1]{Une pyramide à base carrée}
Une pyramide a une base carrée de côté 6 cm et une hauteur de 8 cm. Calculer son volume.
\label{fig-ex-17-2}\figureGeo{sol-02}
\end{exercice}\end{blocLivre}
\begin{blocLivre}{ex-17-3}\begin{exercice}[1]{Un cône}
Un cône a pour rayon 3 cm et pour hauteur 5 cm. Donner son volume exact puis au dixième de cm$^3$.
\label{fig-ex-17-3}\figureGeo{sol-03}
\end{exercice}\end{blocLivre}
\begin{blocLivre}{ex-17-4}\begin{exercice}[1]{Du diamètre au rayon}
Un cône a une base de diamètre 8 cm et une hauteur de 9 cm. Calculer son volume exact.
\label{fig-ex-17-4}\figureGeo{sol-04}
\end{exercice}\end{blocLivre}
\begin{blocLivre}{ex-17-5}\begin{exercice}[1]{Lire trois coordonnées}
Le pavé a pour dimensions 6, 4 et 3 unités, selon les axes $x$, $y$, $z$. Donner les coordonnées du sommet G opposé à O et de H au-dessus du point de coordonnées (0;4;0).
\label{fig-ex-17-5}\figureGeo{sol-05}
\end{exercice}\end{blocLivre}

\rubrique{Entraînement}
\begin{blocLivre}{ex-17-6}\begin{exercice}[2]{Placer un point}
Dans le pavé repéré, placer M(2;3;1). Décrire les trois déplacements depuis O.
\label{fig-ex-17-6}\figureGeo{sol-06}
\end{exercice}\end{blocLivre}
\begin{blocLivre}{ex-17-7}\begin{exercice}[2]{Base triangulaire}
Une pyramide a pour base un triangle de côté 8 cm et de hauteur associée 3 cm. La hauteur de la pyramide est 10 cm. Calculer son volume.
\label{fig-ex-17-7}\figureGeo{sol-07}
\end{exercice}\end{blocLivre}
\begin{blocLivre}{ex-17-8}\begin{exercice}[2]{Retrouver une hauteur}
Une pyramide a un volume de 150 cm$^3$ et une aire de base de 45 cm$^2$. Déterminer sa hauteur.
\label{fig-ex-17-8}\figureGeo{sol-08}
\end{exercice}\end{blocLivre}
\begin{blocLivre}{ex-17-9}\begin{exercice}[2]{Un patron de pyramide}
On veut un patron de pyramide de base carrée de côté 4 cm, dont les quatre arêtes latérales mesurent 5 cm. Décrire puis construire un patron.
\label{fig-ex-17-9}\figureGeo{sol-09}
\end{exercice}\end{blocLivre}
\begin{blocLivre}{ex-17-10}\begin{exercice}[2]{Un patron de cône}
Un cône a pour rayon de base 3 cm et pour génératrice 5 cm. Le représenter en perspective, puis compléter son patron : un disque de base et un secteur de rayon 5 cm dont on déterminera l’angle au centre.
\label{fig-ex-17-10}\figureGeo{sol-10}
\end{exercice}\end{blocLivre}

\rubrique{Problèmes}
\begin{blocLivre}{ex-17-11}\begin{exercice}[3]{Une génératrice calculée}
Un cône a une hauteur de 12 cm et un rayon de 5 cm. Calculer sa génératrice puis son volume exact.
\label{fig-ex-17-11}\figureGeo{sol-11}
\end{exercice}\end{blocLivre}
\begin{blocLivre}{ex-17-12}\begin{exercice}[3]{Comparer un cylindre et un cône}
Un cône et un cylindre ont la même base de rayon 4 cm et la même hauteur de 6 cm. Donner leurs volumes et leur différence.
\label{fig-ex-17-12}\figureGeo{sol-12}
\end{exercice}\end{blocLivre}
\begin{blocLivre}{ex-17-13}\begin{exercice}[3]{Deux unités}
Un réservoir conique a un rayon de 40 cm et une hauteur de 1,2 m. Donner sa capacité en litres au dixième.
\label{fig-ex-17-13}\figureGeo{sol-13}
\end{exercice}\end{blocLivre}
\begin{blocLivre}{ex-17-14}\begin{exercice}[3]{Un assemblage}
Un solide est formé d’un pavé de 8 cm sur 6 cm sur 4 cm, surmonté d’une pyramide dont la base est le rectangle supérieur du pavé et dont la hauteur vaut 3 cm. Calculer le volume total.
\label{fig-ex-17-14}\figureGeo{sol-14}
\end{exercice}\end{blocLivre}
\begin{blocLivre}{ex-17-15}\begin{exercice}[3]{Un remplissage partiel}
Un récipient conique de rayon 6 cm et de hauteur 10 cm est rempli aux trois quarts de sa capacité. Calculer le volume d’eau exact en cm$^3$, puis en mL au dixième. Peut-on en déduire une hauteur d’eau de 7,5 cm ?
\label{fig-ex-17-15}\figureGeo{sol-15}
\end{exercice}\end{blocLivre}

\clearpage
